\documentclass[a4paper,11pt]{scrartcl}

% Language and citation support
\usepackage[english]{babel}
\usepackage[utf8]{inputenc}
\usepackage{csquotes}
\usepackage{placeins}
% Load consolidated formatting from labstyle
\usepackage{labstyle}
\usepackage{cleveref}

% Bibliography setup
\usepackage[
  backend=biber,
  style=numeric,
  sorting=none
]{biblatex}
\addbibresource{Zotero_zitieren.bib}

% PDF inclusion (for title page)
\usepackage{pdfpages}

% Document metadata
\renewcommand{\labtitle}{Drone Control Tuning}
\renewcommand{\labsubtitle}{Development of an Offline Betaflight Drone Control Tuning tool}
\renewcommand{\labauthor}{Yuri Bianchi, Dario Jurietti and Janick Dort}
\renewcommand{\labdate}{\today}

\begin{document}

\begin{titlepage}

\end{titlepage}


\clearpage

\section{Angle Control}

\subsection{Difference between Angle and Acro Mode}
The Betaflight software provides three main flight modes: Angle, Horizon and Acro.

In Angle mode, the drone’s tilt angle is directly linked to the stick position. The maximum tilt is limited to a predefined value, preventing flips and making this mode suitable for beginners.

In contrast, Acro mode provides full manual control. Stick inputs define an angular rate rather than a target angle, meaning the drone continues rotating as long as the stick is deflected. This enables aggressive maneuvers such as flips and rolls but requires greater pilot experience.

During the project thesis, an offline tuner for the angular rate controller was further developed to systematically optimize the rate control loop for improved Acro performance.

The tuner was subsequently extended to support Angle mode, enabling enhanced stability and responsiveness while maintaining intuitive handling for beginner pilots.

\subsection{Control Loop of the Angle Mode}

In Angle mode, the angular rate controller is extended by an outer angle control loop, resulting in a cascaded control structure.

The stick input initially generates a rate setpoint. Unlike in Acro mode, this rate setpoint is converted into a target angle. This target angle then serves as the reference input for the inner angular rate controller.

The cascaded architecture separates attitude control from angular rate stabilization. The outer loop ensures convergence toward the desired orientation, while the inner loop provides fast dynamic stabilization by regulating angular velocity. Consequently, the system achieves improved disturbance rejection and more intuitive handling characteristics, since the pilot directly commands an angle rather than a rotational rate.

The Angle mode controller consists of two cascaded loops, as illustrated in Figure~\ref{fig:angle_control_structure}, and is reconstructed from the Betaflight codebase:

\begin{itemize}
    \item The outer angle loop, which computes the angle error
    \item The inner angular rate loop, which computes the rate error
\end{itemize}

The angle control loop consists of a proportional controller with an additional feedforward term. Within the offline tuning framework, only the proportional controller is optimized. The feedforward term is deactivated and therefore not shown in the figure.

\begin{figure}[h]
    \centering
    \includegraphics[width=1.0\textwidth]{Bilder/Control_Loop_Angle.png}
    \caption{Control loop structure of the Angle control in Betaflight.}
    \label{fig:angle_control_structure}
\end{figure}
\newpage


\subsubsection{Derivation of the Target Angle}

The target angle is obtained by scaling the commanded angular rate with respect to the maximum achievable rate and the configured maximum tilt angle:

\begin{equation}
\theta_{\text{target}} =
\theta_{\text{limit}} \cdot
\frac{\omega_{\text{set}}}{\omega_{\text{max}}}
\end{equation}

where:

\begin{itemize}
    \item $\theta_{\text{limit}}$ is the maximum allowed tilt angle (e.g.\ $60^\circ$),
    \item $\omega_{\text{set}}$ is the commanded angular rate derived from the stick input,
    \item $\omega_{\text{max}}$ is the maximum achievable angular rate determined by the rate configuration.
\end{itemize}

This relation ensures a proportional mapping between stick deflection and target angle.

For example:

\begin{align}
\omega_{\text{set}} &= \omega_{\text{max}}
\Rightarrow
\theta_{\text{target}} = \theta_{\text{limit}}
\\
\omega_{\text{set}} &= 0
\Rightarrow
\theta_{\text{target}} = 0
\end{align}

Thus, the stick position scales the target angle continuously within the range

\begin{equation}
-\theta_{\text{limit}} \leq \theta_{\text{target}} \leq +\theta_{\text{limit}}.
\end{equation}

In Figure x you see the in the upper plot the Gyro signal and in the lower plot the corresponding target angle and angle signal. The target angle is derived from the stick input and the Gyro signal represents the measured angular rate. The angle signal is obtained by integrating the Gyro signal over time, which corresponds to the actual orientation of the drone.

\begin{figure}[h]
    \centering
    \includegraphics[width=1.0\textwidth]{Bilder/Gyro_Angle_Overview.png}
    \caption{Gyro and Angle Signal Overview.}
\end{figure}

\subsection{Chirp Signal for System Identification}

Similar to Acro mode, a chirp signal can also be used for system identification of the Angle control loop. The only difference is that the drone must be operated in Angle mode during excitation.

The chirp signal is added to the current PID setpoint. As a result, the system has two inputs: the stick input, which generates the target angle, and the chirp signal, which is superimposed on the PID setpoint.

\begin{figure}[h]
    \centering
    \includegraphics[width=1.0\textwidth]{Bilder/Chirp_Control_Loop_Angle.png}
    \caption{Gyro signal and corresponding target and measured angle.}
    \label{fig:gyro_angle_overview}
\end{figure}

Thus, the output signal is influenced by both the stick input and the chirp excitation, allowing the system to be excited across a wide frequency range while still maintaining control over the target angle.

However, this also influences the estimation of the transfer function. Therefore, to isolate the plant dynamics, we compute the ratio of the transfer function from target angle to measured angle and the transfer function from target angle to measured angular rate.

The measured angle can be expressed as:

\begin{equation}
\Theta_{\text{measured}}(s)
=
\Theta_{\text{target}}(s) \cdot G_{\text{setpointAngle}}(s)
+
\text{Chirp}(s) \cdot G_{\text{ChirpAngle}}(s)
\end{equation}

The measured angle is therefore composed of two contributions: the closed-loop response to the commanded target angle and the response to the injected chirp excitation. While the stick input defines the desired operating point of the drone, the chirp signal perturbs the system in order to excite a broad frequency range for identification.

To extract the true plant dynamics, the ratio between the transfer function from target angle to measured angle and the transfer function from target angle to measured angular rate is evaluated. In this way, the controller dynamics cancel out and the Angle plant can be isolated.

\newpage

\section{Calculation of the Angle Mode Plant}

For the offline tuner, it is necessary to derive the transfer function of the plant in Angle mode. In a closed-loop control system, the plant represents the physical dynamics that we want to control. In this case, the plant describes the relationship between the commanded angular rate and the resulting angle of the drone.
\subsection{Theoretical Expectation}

If we consider the fundamental relationship between angular rate and orientation,
the angle corresponds to the time integral of the angular rate. Hence, the measured
angle can be written as

\begin{equation}
\theta_{\text{measured}}(t) = \int \omega_{\text{measured}}(t)\, dt.
\end{equation}

Assuming zero initial conditions and applying the Laplace transform, this relationship
in the frequency domain becomes

\begin{equation}
\Theta_{\text{measured}}(s) = \frac{1}{s} \Omega_{\text{measured}}(s).
\end{equation}

This expression shows that the angle is obtained by integrating the angular rate,
which corresponds to a multiplication by $\frac{1}{s}$ in the Laplace domain.
Consequently, from a theoretical point of view, the Angle plant behaves as an ideal integrator with transfer function

\begin{equation}
P_{\text{Angle}}(s) = \frac{1}{s}.
\end{equation}

\subsection{Closed-Loop Identification Problem}

However, during flight the system operates in closed loop, meaning that the measured response is influenced by the control actions of the loop. Therefore, the input signal is no longer independent of the output.

If we were to directly estimate the transfer function from input to output,

\begin{equation}
P(s)_{\text{Angle}} = \frac{\Theta_{\text{measured}}(s)}{\Omega_{\text{measured}}(s)}
\end{equation}

the result would not reflect the true plant dynamics. Instead, it would include:

\begin{itemize}
    \item the dynamics of the controller,
    \item the feedback effects of the closed loop,
    \item and the plant dynamics.
\end{itemize}

Therefore, the estimated transfer function would deviate significantly from the theoretical expectation of an integrator and would not be suitable for tuning the controller parameters.

\subsection{Closed-Loop Transfer Relations}

To overcome this issue, we can use the known relationships between the transfer functions in a closed-loop system. 

We can calculate the transfer function of the complete system from the target angle to the measured angle, which is not directly affected by the closed-loop feedback:

\begin{equation}
G_{\text{setpoint->Current Angle}}= \frac{\Theta_{\text{measured}}(s)}{\Theta_{\text{target}}(s)}
\end{equation}
\begin{equation}
= \frac{
P_G(s)\,
P_A(s)\,
C_A(s)\,
C_{\text{PI},G}(s)
}{
1
+ P_G(s)\,
C_{\text{PI},G}(s)\,
C_{\text{D},G}(s)
+ P_A(s)\,
P_G(s)\,
C_A(s)\,
C_{\text{PI},G}(s)
}
\end{equation}

\noindent
\textbf{Legend:}
\begin{itemize}
    \item $P_G(s)$: Angular rate plant
    \item $P_A(s)$: Angle plant
    \item $C_A(s)$: Angle controller
    \item $C_{\text{PI},G}(s)$: PI controller of the angular rate loop
    \item $C_{\text{D},G}(s)$: D controller of the angular rate loop
\end{itemize}
\par

Additional to this, we also have to calculate the transfer function from the target angle to the measured angular rate:

\begin{equation}
G_{\text{setpoint->Current Angular Rate}}(s) = \frac{\Omega_{\text{measured}}(s)}{\Theta_{\text{target}}(s)} 
\end{equation}
\begin{equation}
= \frac{
P_G(s)\,
C_A(s)\,
C_{\text{PI},G}(s)
}{
1
+ P_G(s)\,
C_{\text{PI},G}(s)\,
C_{\text{D},G}(s)
+ P_A(s)\,
P_G(s)\,
C_A(s)\,
C_{\text{PI},G}(s)
}
\end{equation}

With these two transfer functions, we can derive the Angle plant by taking the ratio of the transfer function from target angle to measured angle and the transfer function from target angle to measured angular rate.
Since both expressions share the same denominator and contain identical controller dynamics, all common terms cancel out. As a result, only the Angle plant $P_A(s)$ remains. This allows us to isolate the true plant dynamics of the Angle mode independently of the controller structure.

\begin{equation}
P_{\text{Angle}}(s) =
\frac{
G_{\text{setpoint->Current Angle}}(s)
}{
G_{\text{setpoint->Current Angular Rate}}(s)
}
\end{equation}

This approach allows us to isolate the true plant dynamics of the Angle mode, enabling accurate tuning of the controller parameters based on the estimated transfer function.

\newpage


\subsection{Results of the System Identification}

Through the application of the chirp signal and the estimation of the transfer functions, as mentioned in the previous section, we can derive the estimated transfer function of the Angle plant. The resulting Bode plot of the estimated transfer function is shown in Figure~\ref{fig:angle_bode}.

\begin{figure}[h]
    \centering
    \includegraphics[width=1.0\textwidth]{Bilder/Angle_Plant_Bode.png}
    \caption{Bode plot of the estimated Angle plant.}
    \label{fig:angle_bode}
\end{figure}

As you can see, the estimated transfer function closely follows the theoretical expectation of an integrator at lower frequencies, confirming that the Angle plant behaves as expected. However, at higher frequencies, deviations from the ideal integrator behavior can be observed. These deviations can be attributed to various factors such as sensor noise or unmodeled dynamics in the system.

The coherence in the plot indicates the reliability of the estimation across the frequency range. A high coherence value, especially at lower frequencies, suggests that the estimated transfer function accurately captures the system dynamics in that range. Conversely, a drop in coherence at higher frequencies indicates that the estimation may be less reliable in those regions, likely due to noise or other disturbances affecting the measurements.

To obtain the coherence value, the following formula can be used:

\begin{equation}
\gamma^2(f) = \frac{|P_{xy}(f)|^2}{P_{xx}(f)\, P_{yy}(f)}
\end{equation}

where:
\begin{itemize}
    \item $\gamma^2(f)$ is the coherence at frequency $f$,
    \item $P_{xy}(f)$ is the cross power spectral density between the input and output signals,
    \item $P_{xx}(f)$ is the power spectral density of the input signal,
    \item $P_{yy}(f)$ is the power spectral density of the output signal.
\end{itemize}



\end{document}
